\documentclass[11pt]{article}

\usepackage{notes}        % custom package (optional)
\usepackage{amsmath}
\usepackage{amssymb}
\usepackage{geometry}
\usepackage{graphicx}
\usepackage{enumitem}

\geometry{margin=1in}

\title{Fourier Transforms with Potential Fields \\
Instructor Key and Interpretation Guide}
\author{Geophysics IIIA \\ D. Hasterok}
\date{}

\begin{document}
\maketitle

\section*{Purpose of This Document}

This document provides an instructor-only interpretation guide and marking key for
the Jupyter notebook practical on Fourier transforms applied to gravity and magnetic
data from South Australia.  It is intended to clarify expected observations, physical
interpretations, and common student misconceptions.  Students are not expected to
provide all of the detail included here; rather, this guide defines what constitutes
a sound geophysical understanding for each question.

\section*{Notebook Structure}

The practical notebook (\texttt{week\_7\_fourier.ipynb}) is organised into six
sections:
\begin{enumerate}
  \item Load and display gravity, magnetics, and topography/sediment grids
  \item Interactive profile selection (click on a map, or use fallback coordinates)
  \item One-dimensional profile continuation with an interactive slider
  \item Two-dimensional gravity continuation and isostatic comparison
  \item Reduction to the pole (interactive inclination and declination sliders)
  \item Sediment-thickness correction for RTP (\emph{optional/exploratory})
\end{enumerate}
The problem set (\texttt{week7\_problem\_set.tex}) mirrors this structure with 20
questions grouped into five assessed sections plus one optional section.

\section*{Learning Outcomes}

By the end of this practical, students should be able to:
\begin{itemize}
  \item Interpret spatial and spectral representations of potential-field data
  \item Explain upward and downward continuation as frequency-domain operations
  \item Distinguish regional and residual fields and critically evaluate the
        choices involved
  \item Compare spectral and isostatic approaches to defining a regional gravity
        field
  \item Apply reduction to the pole and understand its assumptions and limitations
  \item Discuss how sediment thickness affects the interpretation of magnetic data
  \item Appreciate the non-uniqueness and ill-posed nature of potential-field
        interpretation
\end{itemize}

% ─────────────────────────────────────────────────────────────────────────────
\section{Section~1 — One-Dimensional Profile Continuation (Q1--5)}
% ─────────────────────────────────────────────────────────────────────────────

Students extract co-located gravity and magnetic profiles along a user-selected
great-circle path, then apply the 1-D continuation filter
$H(k, z) = e^{-2\pi|k|z}$ (wavenumber in cycles~km$^{-1}$, height $z$ in km)
using an interactive slider from $-20$ to $+100$~km.

\subsection*{Q1 — Upward continuation as a low-pass filter}

\paragraph{Expected answer.}
Narrow, localised anomalies disappear first because they correspond to high
spatial wavenumbers; the exponential decay factor $e^{-2\pi|k|z}$ is smallest
for large $|k|$.  Broad, regional features persist longest because they are
associated with small $|k|$.  This implies that short-wavelength anomalies
originate from shallow sources, while long-wavelength anomalies reflect deeper or
more laterally extensive structure.

\paragraph{Common misconceptions.}
\begin{itemize}
  \item Claiming that broad anomalies disappear first.
  \item Confusing amplitude attenuation with frequency-independent smoothing.
\end{itemize}

\subsection*{Q2 — Residual anomalies}

\paragraph{Expected answer.}
The residual (observed $-$ regional) emphasises sources not accounted for by the
long-wavelength regional, typically shallower lithological contrasts, intrusive
bodies, or structural offsets.  Gravity and magnetic residuals are often
spatially correlated where lithology governs both density and susceptibility
(e.g.\ mafic intrusions), but may disagree where remanent magnetisation decouples
the two fields, or where a density contrast exists without a corresponding
susceptibility contrast.

\subsection*{Q3 — Choosing a continuation height}

\paragraph{Expected answer.}
\begin{enumerate}[label=\alph*.]
  \item Defensible criteria include: visual inspection (the height at which known
        deep regional features become smooth while local anomalies remain visible
        in the residual); comparison with independent depth estimates (e.g.\ Moho
        depth from seismic); spectral depth analysis (slope of the log-power
        spectrum); or geological judgement about the expected crustal structure.
  \item Sensitivity is typically high.  Students should observe that the
        character of the residual changes substantially over a range of only a
        few kilometres, which is a key point about the subjective nature of the
        regional--residual separation.
\end{enumerate}

\subsection*{Q4 — Instability of downward continuation}

\paragraph{Expected answer.}
\begin{enumerate}[label=\alph*.]
  \item Visible instability (ringing, oscillations, extreme amplitudes) typically
        sets in within a few kilometres of downward continuation.  The exact
        depth depends on the noise level of the dataset.
  \item Magnetic data tend to become unstable at shallower depths than gravity
        because magnetic datasets typically have higher spatial frequency content
        and more noise.
  \item The filter $e^{+2\pi|k||z|}$ amplifies all wavenumbers; high-wavenumber
        noise grows fastest.  The effective maximum stable depth is set by the
        noise floor of the measurement, which differs between the two datasets.
\end{enumerate}

\subsection*{Q5 — Physical meaning of the continued field}

\paragraph{Expected answer.}
Yes, the upward-continued profile should look physically reasonable: it represents
the field that would be measured on a higher observation surface.  It should be
smoother and lower in amplitude.  Features that seem unreasonable (e.g.\
negative-amplitude oscillations flanking a positive body) may reflect edge effects
from the FFT (wrap-around), or the presence of sources that violate the
harmonic-function assumption (e.g.\ sources within the continuation region).

% ─────────────────────────────────────────────────────────────────────────────
\section{Section~2 — 2-D Gravity Continuation and Isostasy (Q6--9)}
% ─────────────────────────────────────────────────────────────────────────────

The 2-D continuation filter $H = e^{-2\pi\sqrt{k_x^2+k_y^2}\,z}$ is applied to
the full gravity grid with a fixed reference height of 25~km for the comparison
maps.  An Airy isostatic correction is computed as
$g_\mathrm{iso} = 2\pi G\,\rho_c\, h$, where $h$ is positive topography
(sub-sea-level values are zeroed).

\subsection*{Q6 — Spectral regional vs isostatic correction}

\paragraph{Expected answer.}
The two fields share long-wavelength trends but differ in detail.
\begin{enumerate}[label=\alph*.]
  \item Broad agreement is expected in elevated terrain where Airy compensation
        is a reasonable approximation.
  \item The largest differences occur in areas of thick sedimentary basins (which
        have near-zero topography but significant subsurface density variation),
        along the coast (ocean loading is not captured), and in regions of
        tectonic flexure where isostasy is not locally achieved.
\end{enumerate}

\subsection*{Q7 — Residual comparison and compensation depth}

\paragraph{Expected answer.}
Students should find rough visual agreement at continuation heights in the range
of 30--50~km, consistent with compensation depths for cratonic crust.  This does
not prove compensation depth equals that height; it reflects the scale of the
features that Airy compensation removes.

\subsection*{Q8 — Anomalies not explained by isostasy}

\paragraph{Expected answer.}
Anomalies in the spectral residual but absent from the isostatic residual reflect
features accounted for by upward continuation but not by Airy topographic
compensation, e.g.\ deep crustal density contrasts, underplated mafic material,
or lateral mantle heterogeneity.  The converse (present in isostatic but not
spectral residual) may indicate shallow features whose gravity is removed by
continuation but not by the topographic model.

\subsection*{Q9 — Limitations of the Airy model}

\paragraph{Expected answer.}
Key assumptions and their failure modes:
\begin{itemize}
  \item \emph{Local compensation:} flexural rigidity of the lithosphere means
        loads are regionally compensated; the Airy model therefore over-corrects
        beneath narrow ranges and under-corrects adjacent to them.
  \item \emph{Uniform crustal density $\rho_c$:} lateral variations in crustal
        composition are ignored.
  \item \emph{Topography as the only load:} sediment thickness, ice loading, and
        internal density anomalies are not accounted for.
  \item \emph{Flat Moho geometry:} the model assumes a planar reference Moho,
        ignoring pre-existing crustal structure.
\end{itemize}

% ─────────────────────────────────────────────────────────────────────────────
\section{Section~3 — Reduction to the Pole (Q10--13)}
% ─────────────────────────────────────────────────────────────────────────────

The stabilised RTP filter
$H_\mathrm{RTP} = \bar{\Theta}^2 / (|\Theta|^2 + \varepsilon)^2$
is applied to the total magnetic intensity (TMI) grid.  Physical spacing in km
is used for the wavenumber axes.  Students explore inclination ($-90°$ to $-10°$)
and declination ($-30°$ to $+30°$) interactively.

\subsection*{Q10 — Effect on anomaly positions}

\paragraph{Expected answer.}
In the Southern Hemisphere, positive TMI anomalies are displaced southward from
their sources due to the negative inclination of the ambient field.  RTP shifts
anomaly maxima northward, centring them above the causative bodies.  The shift
magnitude depends on inclination; at $I \approx -60°$ (typical for South
Australia) the displacement is moderate, of order 5--20~km for typical anomaly
widths.  Shape changes from asymmetric (TMI) to more symmetric (RTP) should also
be noted.

\subsection*{Q11 — Stability at low inclinations}

\paragraph{Expected answer.}
\begin{enumerate}[label=\alph*.]
  \item Artefacts typically appear at inclinations shallower than about $\pm15°$
        to $\pm20°$, though the water-level parameter $\varepsilon = 0.05$
        provides some stability.  Students should see ringing and extreme
        amplitudes in the RTP map as the slider approaches $0°$.
  \item Near the equator ($I \approx 0$), the direction-cosine factor $\Theta$
        approaches zero for wavenumbers in the north--south direction, making the
        filter singular.  The inversion is numerically unstable even with
        regularisation, so RTP cannot be reliably applied at low magnetic
        latitudes.
\end{enumerate}

\subsection*{Q12 — Difference map (RTP $-$ TMI)}

\paragraph{Expected answer.}
The difference map highlights the phase correction — positive differences appear
to the north of anomaly centres (where RTP has moved the peak) and negative
differences to the south.  This dipolar difference pattern encodes the along-field
asymmetry of the original TMI anomaly.  Larger differences indicate either a
stronger phase shift (lower inclination, shallower source) or stronger
magnetisation contrast.

\subsection*{Q13 — Remanent magnetisation}

\paragraph{Expected answer.}
The RTP filter applies a single phase correction based on the ambient-field
direction.  Where a body carries remanent magnetisation in a significantly
different direction (e.g.\ a reversely magnetised unit), RTP will
\emph{over-correct} or \emph{under-correct} the anomaly, potentially distorting
its shape and position rather than centring it.  Students who have covered
remanence (e.g.\ in Week 6) may identify candidate anomalies as those that do
not become symmetric after RTP, or that show unexpected polarity.

% ─────────────────────────────────────────────────────────────────────────────
\section{Section~4 — Sediment-Thickness Correction (Optional, Q14--17)}
% ─────────────────────────────────────────────────────────────────────────────

This section is explicitly marked optional in the problem set.  It is intended as
an exploratory exercise.  The combined filter applies RTP plus a downward
continuation of $z_\mathrm{sed}$~km damped by a Gaussian taper of length scale
$\lambda$~km.

\subsection*{Q14 — Where the correction matters most}

\paragraph{Expected answer.}
The largest differences between standard and sediment-corrected RTP should appear
in regions of thickest sediment (the Officer Basin, Murray Basin, and Eucla
Basin).  Students should verify this against the sediment-thickness map shown in
Section~1 of the notebook.

\subsection*{Q15 — Role of the damping parameter $\lambda$}

\paragraph{Expected answer.}
Increasing $\lambda$ (the Gaussian taper length scale) progressively damps the
short-wavelength amplification of the sediment correction.  At large $\lambda$,
the correction effectively removes the sediment-depth inversion entirely for
high wavenumbers, leaving only the long-wavelength component.  Whether this is
geologically sensible depends on the data noise level; students should observe
that large $\lambda$ produces a smooth but geologically uninformative correction,
while small $\lambda$ produces artefacts.

\subsection*{Q16 — Validity of the uniform approximation}

\paragraph{Expected answer.}
\begin{enumerate}[label=\alph*.]
  \item The approximation is most valid where sediment thickness is close to the
        mean value, i.e.\ in moderately covered areas away from basin margins.
        It is least valid in the centres of thick basins (over-correction) and
        in exposed basement terranes (unnecessary correction).
  \item A rigorous spatially variable correction would require either a
        depth-to-basement model at grid resolution, a tiled spectral approach
        (windowed FFT with spatial overlap), or iterative modelling.
\end{enumerate}

\subsection*{Q17 — Analogous gravity correction}

\paragraph{Expected answer.}
A downward-continuation gravity correction for sediment would attempt to recover
the gravity field at the basement surface rather than the observation surface.
It would differ from the Airy correction in that: (1) the Airy correction removes
the long-wavelength gravity effect of the isostatic root driven by topography,
whereas a sediment downward continuation attempts to resolve the gravity field at
depth; and (2) a downward continuation is inherently unstable, while the Airy
correction is a stable forward model.  Both approaches are geologically
approximate.

% ─────────────────────────────────────────────────────────────────────────────
\section{Section~5 — Synthesis (Q18--20)}
% ─────────────────────────────────────────────────────────────────────────────

\subsection*{Q18 — Wavelength and depth}

\paragraph{High-quality answer.}
The approximate relationship is $\lambda \approx 2\pi\,d$, where $\lambda$ is the
dominant wavelength and $d$ is the source depth (derived from the radial decay of
a point source).  Key assumptions are: isolated, compact sources; source-free
continuation region; and stationary noise statistics.  A misleading depth estimate
arises when a large, laterally extensive source (e.g.\ a thick sedimentary basin)
dominates the long-wavelength signal — the ``depth'' inferred from the spectral
slope would reflect the lateral extent of the source, not its burial depth.

\subsection*{Q19 — Spectral vs isostatic regional field}

\paragraph{High-quality answer.}
\emph{Spectral continuation} is mathematically objective and requires no
geological model, but the choice of height is arbitrary and geologically
unmotivated; it works well for separating sources at different depths.
\emph{Isostatic correction} is physically motivated and removes the predicted
gravity of crustal roots, but assumes a specific compensation mechanism (Airy),
requires density parameters, and fails in regions of flexural or dynamic
topography.  The spectral approach is preferred for exploratory analysis; the
isostatic approach is preferred where the geological setting supports local Airy
compensation and the goal is to isolate intra-crustal anomalies.

\subsection*{Q20 — Non-uniqueness}

\paragraph{High-quality answer.}
A concrete example: at a suitable continuation height, the spectral residual
gravity map closely resembles the isostatic residual gravity map, yet the two
are derived from entirely different assumptions (mathematical filter vs.\
physical compensation model).  Both ``explain'' the observations, but they imply
different geological models.  Discriminating between them would require
independent constraints such as seismic refraction or reflection profiles (to
determine actual Moho depth and crustal thickness), well data (to constrain
density and susceptibility), or receiver function analysis.

% ─────────────────────────────────────────────────────────────────────────────
\section*{Marking Guidance}
% ─────────────────────────────────────────────────────────────────────────────

\begin{itemize}
  \item \textbf{Full marks} require correct physical reasoning, not just
        description of what is seen on screen.
  \item \textbf{Partial marks} for correct observation without physical
        explanation, or correct explanation applied to an incorrect observation.
  \item \textbf{No marks} for purely descriptive statements that restate the
        question (e.g.\ ``the continued field is smoother because continuation
        makes it smoother'').
  \item The synthesis questions (Q18--20) are worth more and should demonstrate
        integration across sections.  Accept well-reasoned answers that differ
        from the expected answer provided the geophysical logic is sound.
  \item The optional sediment correction section (Q14--17) is marked generously
        given its exploratory nature; acknowledge any thoughtful engagement with
        the instability and approximation issues.
\end{itemize}

\vspace{1em}
\noindent\textbf{End of Instructor Key}

\end{document}
